%============================================================
% CH1 · P1 · PROBLEM 2 — ADVANTAGE (Codeforces 1760C)
%============================================================
\subsection{Problem 1.1.2: Advantage}
\problemlabel{1.1.2}{ssec:ch1-advantage}

\problemcard{1.1.2}{Advantage}{Codeforces 1760C}{https://codeforces.com/problemset/problem/1760/C}{\exeasy\ \cf{800}}{Two largest values (duplicates allowed), then a second pass}

\subsubsection*{1. Problem statement}

There are $n$ participants; participant $i$ has strength $s_i$. Each participant wants to know how much stronger (or weaker) he is than the \textbf{strongest of the other participants}. For every $i$, print
\[
d_i \;=\; s_i - \max_{j \ne i} s_j ,
\]
which may be negative. The input contains $t$ independent test cases.

\textbf{Input:} the first line contains $t$; each test case consists of a line with $n$ and a line with $s_1, \dots, s_n$. \textbf{Output:} for each test case, the $n$ numbers $d_1, \dots, d_n$ on one line, separated by spaces.

\textbf{Constraints:} $1 \le t \le 1000$, \; $2 \le n \le 2\cdot10^5$, \; $1 \le s_i \le 10^9$, \; the sum of $n$ over all test cases does not exceed $2\cdot10^5$.

\begin{table}[H]
\centering\small
\begin{tabular}{@{}lll@{}}
\toprule
\textbf{test} & \textbf{strengths} & \textbf{output} \\
\midrule
1 & $4\ 7\ 3\ 5$ & $-3\ 2\ -4\ -2$ \\
2 & $1\ 2$ & $-1\ 1$ \\
3 & $1\ 2\ 3\ 4\ 5$ & $-4\ -3\ -2\ -1\ 1$ \\
4 & $4\ 9\ 4$ & $-5\ 5\ -5$ \\
5 & $4\ 4\ 4\ 4$ & $0\ 0\ 0\ 0$ \\
\bottomrule
\end{tabular}
\end{table}

\begin{diagrambox}
\begin{tikzpicture}
    \begin{scope}
        \dparray{s}{0}{$s$}{4/acell, 7/aread, 3/acell, 5/akeep}
        \dparray{d}{-0.85}{$d$}{$-3$/compcell, 2/anscell, $-4$/compcell, $-2$/compcell}
        \node[note, anchor=north, text width=4.4cm] at (1.8,-1.4) {test 1: everybody except participant 2 compares with $7$; participant 2 compares with $5$};
    \end{scope}
    \begin{scope}[xshift=6.6cm]
        \dparray{t}{0}{$s$}{4/akeep, 4/akeep, 4/akeep, 4/akeep}
        \dparray{e}{-0.85}{$d$}{0/compcell, 0/compcell, 0/compcell, 0/compcell}
        \node[note, anchor=north, text width=4.4cm] at (1.8,-1.4) {test 5: the strongest \emph{other} participant always has strength $4$};
    \end{scope}
\end{tikzpicture}
\figcap{Official tests 1 and 5. Orange: the participant whose opponent is the second strongest; green: the strongest other participant.}
\end{diagrambox}

\subsubsection*{2. The problem model}
\begin{modelbox}[Problem Model --- Advantage]
\renewcommand{\arraystretch}{1.35}
\begin{tabularx}{\linewidth}{@{}>{\bfseries}l X@{}}
Input & $t$ arrays of positive integers; total length at most $2\cdot10^5$ \\
Computed & an \emph{array}: for every position $i$, $s_i$ minus the maximum of the array with position $i$ removed \\
Over & elements (positions count separately: two equal strengths are two participants) \\
Sizes & $n \le 2\cdot10^5$ per test and in total: $\bigO{n}$ or $\bigO{n \log n}$; $\bigO{n^2} = 4\cdot10^{10}$ is far too slow \\
Values & $|d_i| < 10^9$: fits in \texttt{int}; no sums are needed \\
\end{tabularx}
\end{modelbox}

\begin{diagrambox}
\begin{tikzpicture}
    \dparray{m}{0}{$s$}{3/acell, 8/akeep, 2/acell, 8/akeep, 5/acell}
    \node[font=\scriptsize\bfseries, text=deepblue, anchor=west] at (5.1,0.45) {global summary};
    \node[rectangle, rounded corners=3pt, draw=successgreen, fill=solutionbg, font=\scriptsize, align=left,
          anchor=west, text width=3.6cm] at (5.1,-0.15) {$m_1 = 8$, \; $m_2 = 8$\\(two largest \emph{elements})};
    \node[note, anchor=north west, text width=8.8cm, align=left] at (0,-0.75)
        {$\max_{j \ne i} s_j = m_1$ unless $s_i$ \emph{is} the maximum; then it is $m_2$ --- which equals $m_1$ when the maximum occurs twice. So $d = (-5, 0, -6, 0, -3)$.};
\end{tikzpicture}
\figcap{The whole problem reduces to two numbers computed once per test: the largest and the second largest element, duplicates included.}
\end{diagrambox}

\subsubsection*{3. Key observations}

\begin{itemize}[leftmargin=*]
    \item \textbf{Removing one element changes the maximum only if it was the maximum.} For $s_i < m_1$ the strongest other participant is still the overall maximum $m_1$. For $s_i = m_1$ it is the largest of the rest, which is $m_2$, the second largest \emph{element}.
    \item \textbf{Subtracting the maximum from everybody is wrong.} It gives $d = (-3, 0, -4, -2)$ on test 1: the strongest participant gets $0$ instead of $7 - 5 = 2$.
    \item \textbf{The second maximum must count duplicates.} With the second largest \emph{distinct} value, the array $(3, 8, 2, 8, 5)$ gives each 8 the answer $8 - 5 = 3$, but each 8 faces the other 8: the answer is $0$. The same mistake breaks test 5.
    \item \textbf{One pass cannot answer early positions.} When we stand at $i = 0$ we have not seen the rest yet, so the answers need a second pass after the global summary --- the ``two passes'' scheme of Lesson~\ref{ssec:ch1-scanning}.
\end{itemize}

\begin{diagrambox}
\begin{tikzpicture}
    \begin{scope}
        \dparray{w}{0}{$s$}{3/acell, 8/aread, 2/acell, 8/aread, 5/acell}
        \dparray{x}{-0.85}{wrong}{$-5$/compcell, 3/aread, $-6$/compcell, 3/aread, $-3$/compcell}
        \dparray{y}{-1.7}{right}{$-5$/compcell, 0/akeep, $-6$/compcell, 0/akeep, $-3$/compcell}
    \end{scope}
    \node[note, anchor=west, text width=4.6cm, align=left] at (5.2,-0.85) {``wrong'' uses the second largest \emph{distinct} value ($5$); each 8 should be compared with the other 8};
\end{tikzpicture}
\figcap{The counter-example for the distinct second maximum (outputs of both versions).}
\end{diagrambox}

\subsubsection*{4. Approaches}

\begin{bruteforcebox}
\textbf{Idea.} For every $i$, scan all $j \ne i$ and take the maximum.
\begin{lstlisting}[style=cpp]
#include <bits/stdc++.h>
using namespace std;

int main() {
    int t;
    cin >> t;
    while (t--) {
        int n;
        cin >> n;
        vector<int> s(n);
        for (int& x : s) cin >> x;
        for (int i = 0; i < n; i++) {
            int best = 0;                                // strengths are >= 1
            for (int j = 0; j < n; j++)
                if (j != i) best = max(best, s[j]);
            cout << s[i] - best << (i + 1 < n ? ' ' : '\n');
        }
    }
}
\end{lstlisting}
\begin{complexitybox}
\textbf{Time:} $\bigO{n^2}$ per test --- $(2\cdot10^5)^2 = 4\cdot10^{10}$ steps for one big test: \tle{}. It recomputes almost the same maximum $n$ times. \textbf{Space:} $\bigO{n}$. We keep it as the reference for stress testing.
\end{complexitybox}
\end{bruteforcebox}

\begin{algorithmbox}
\textbf{Idea.} Pass 1 keeps the two largest elements $m_1 \ge m_2$ (the non-distinct top-2 of Pattern 1: a new value larger than $m_1$ demotes $m_1$ to $m_2$; a value equal to $m_1$ falls through to the second test and becomes $m_2$). Pass 2 prints $s_i - m_2$ if $s_i = m_1$, else $s_i - m_1$.
\begin{lstlisting}[style=cpp]
#include <bits/stdc++.h>
using namespace std;

int main() {
    ios::sync_with_stdio(false);
    cin.tie(nullptr);
    int t;
    cin >> t;
    while (t--) {
        int n;
        cin >> n;
        vector<int> s(n);
        for (int& x : s) cin >> x;
        int mx1 = 0, mx2 = 0;                          // two largest, duplicates allowed
        for (int x : s) {                              // pass 1
            if (x > mx1) { mx2 = mx1; mx1 = x; }       // new maximum: old one becomes second
            else if (x > mx2) mx2 = x;                 // x == mx1 lands here: second = mx1
        }
        for (int i = 0; i < n; i++) {                  // pass 2
            int best = (s[i] == mx1) ? mx2 : mx1;      // strongest OTHER participant
            cout << s[i] - best << (i + 1 < n ? ' ' : '\n');
        }
    }
}
\end{lstlisting}
\begin{complexitybox}
\textbf{Time:} $\bigO{n}$ per test, $\bigO{\sum n} = \bigO{2\cdot10^5}$ in total. \textbf{Space:} $\bigO{n}$ for the strengths (we must remember them to print the answers after pass 1). \textbf{Verdict:} \ac{}.
\end{complexitybox}
\end{algorithmbox}

Starting both maxima at $0$ is safe here because every strength is at least 1 and $n \ge 2$, so after pass 1 both are real strengths. (On a problem with negative values, start from \texttt{INT\_MIN} or from the first two elements.)

\begin{invariantbox}[Why the answer is correct]
\textbf{Pass 1 invariant.} After reading $s_1, \dots, s_k$, \texttt{mx1} is the largest and \texttt{mx2} the second largest element of this prefix as a \emph{multiset} (with $0$ for a missing one). For the step with value $x$: the two largest of the prefix plus $x$ are the two largest of $\{\texttt{mx1}, \texttt{mx2}, x\}$; if $x > \texttt{mx1}$ they are $(x, \texttt{mx1})$; if $\texttt{mx2} < x \le \texttt{mx1}$ they are $(\texttt{mx1}, x)$ --- including $x = \texttt{mx1}$, a second copy of the maximum; otherwise they do not change. After the pass $m_1 \ge m_2$ are the two largest elements of the array, and $n \ge 2$ makes both real.

\textbf{Pass 2.} Remove position $i$ from the multiset. If $s_i < m_1$, the element $m_1$ (at some other position) remains and is still the maximum. If $s_i = m_1$, one copy of $m_1$ disappears and the maximum of the rest is the second largest element $m_2$ (equal to $m_1$ if the maximum occurred at least twice). In both cases \texttt{best} $= \max_{j \ne i} s_j$. \hfill$\blacksquare$
\end{invariantbox}

\begin{alternativebox}
\textbf{Prefix and suffix maxima.} Without any case analysis: let $\mathit{pre}[i] = \max(s_0, \dots, s_{i-1})$ and $\mathit{suf}[i] = \max(s_i, \dots, s_{n-1})$ (with $0$ for an empty side). The strongest participant other than $i$ is $\max(\mathit{pre}[i], \mathit{suf}[i+1])$ --- the best on the left combined with the best on the right. This is the ``left side and right side'' situation of Lesson~\ref{ssec:ch1-scanning} and works for any operation that can be combined (minimum, sum, OR, \ldots), not only for the maximum.
\begin{lstlisting}[style=cpp]
#include <bits/stdc++.h>
using namespace std;

int main() {
    ios::sync_with_stdio(false);
    cin.tie(nullptr);
    int t;
    cin >> t;
    while (t--) {
        int n;
        cin >> n;
        vector<int> s(n);
        for (int& x : s) cin >> x;
        // pre[i] = max of s[0..i), suf[i] = max of s[i..n); 0 stands for an empty side
        vector<int> pre(n + 1, 0), suf(n + 1, 0);
        for (int i = 0; i < n; i++) pre[i + 1] = max(pre[i], s[i]);
        for (int i = n - 1; i >= 0; i--) suf[i] = max(suf[i + 1], s[i]);
        for (int i = 0; i < n; i++)
            cout << s[i] - max(pre[i], suf[i + 1]) << (i + 1 < n ? ' ' : '\n');
    }
}
\end{lstlisting}
Same $\bigO{n}$ time, $\bigO{n}$ extra memory for the two arrays. It is the method to remember when the operation has no ``second best'' shortcut --- for example ``the GCD of all the others'' (a GCD cannot be undone, but prefix and suffix GCDs combine).

\begin{diagrambox}
\begin{tikzpicture}
    \dpindex{ai}{1.3}{$i$}{0}{4}
    \dparray{as}{0.6}{$s$}{4/acell, 7/aread, 3/acell, 5/acell, {}/idx}
    \dparray{ap}{-0.25}{\textit{pre}}{0/basecell, 4/compcell, 7/compcell, 7/compcell, 7/compcell}
    \dparray{au}{-1.1}{\textit{suf}}{7/compcell, 7/compcell, 5/compcell, 5/compcell, 0/basecell}
    \draw[accentorange, line width=1.2pt, rounded corners=2pt] ([xshift=-1pt,yshift=1pt]ap1.north west) rectangle ([xshift=1pt,yshift=-1pt]ap1.south east);
    \draw[accentorange, line width=1.2pt, rounded corners=2pt] ([xshift=-1pt,yshift=1pt]au2.north west) rectangle ([xshift=1pt,yshift=-1pt]au2.south east);
    \node[note, anchor=west, text width=4.3cm, align=left] at (5.0,-0.25) {for $i = 1$: $\max(\mathit{pre}[1], \mathit{suf}[2]) = \max(4, 5) = 5$, so $d_1 = 7 - 5 = 2$};
\end{tikzpicture}
\figcap{Test 1: the two framed cells are the best participant left of $i = 1$ and the best right of it. Both arrays have $n + 1$ cells; the grey ones stand for an empty side.}
\end{diagrambox}
\end{alternativebox}

\begin{tricksbox}
\begin{itemize}[leftmargin=*]
    \item \textbf{Sorting a copy} gives $m_1$ and $m_2$ as its last two elements: $\bigO{n \log n}$, also \ac{}, and impossible to get the duplicate case wrong.
    \item \textbf{Count of the maximum.} Equivalent formulation: keep $m_1$, its count $c$, and the largest value below $m_1$. Then $\max_{j\ne i} s_j = m_1$ unless $s_i = m_1$ and $c = 1$. More state, same result --- the two-largest version is shorter.
    \item \textbf{Output.} Up to $2\cdot10^5$ numbers: use \texttt{'\textbackslash n'} and the two fast-I/O lines of Lesson~\ref{ssec:ch1-cost}; printing with \texttt{endl} after every number would flush $2\cdot10^5$ times.
    \item \textbf{No per-test reset} is needed beyond the local vector: all work is proportional to this test's $n$.
\end{itemize}
\end{tricksbox}

\subsubsection*{5. Trace}

\begin{tracebox}[{Pass 1 on test 1, $s = (4, 7, 3, 5)$}]
\centering
\begin{tikzpicture}
    \begin{scope}
        \dparray{p}{0}{after $7$}{4/adone, 7/aread, 3/acell, 5/acell}
        \node[note, anchor=west, text width=5cm, align=left] at (3.9,0) {$7 > \texttt{mx1} = 4$: $(\texttt{mx1}, \texttt{mx2}) = (7, 4)$};
    \end{scope}
    \begin{scope}[yshift=-1.0cm]
        \dparray{q}{0}{after $5$}{4/adone, 7/adone, 3/adone, 5/aread}
        \node[note, anchor=west, text width=5cm, align=left] at (3.9,0) {$4 < 5 \le 7$: $\texttt{mx2} = 5$; the strongest participant will be compared with $5$};
    \end{scope}
\end{tikzpicture}

\medskip
\small
\begin{tabular}{@{}cccc@{}}
\toprule
\textbf{read} $x$ & \textbf{case} & \texttt{mx1} & \texttt{mx2} \\
\midrule
(start) & & 0 & 0 \\
4 & $x > \texttt{mx1}$ & 4 & 0 \\
7 & $x > \texttt{mx1}$ & 7 & 4 \\
3 & no change & 7 & 4 \\
5 & $\texttt{mx2} < x \le \texttt{mx1}$ & 7 & 5 \\
\bottomrule
\end{tabular}
\qquad
\begin{tabular}{@{}cccc@{}}
\toprule
$i$ & $s_i$ & \texttt{best} & $d_i$ \\
\midrule
1 & 4 & 7 & $-3$ \\
2 & 7 & 5 & $2$ \\
3 & 3 & 7 & $-4$ \\
4 & 5 & 7 & $-2$ \\
\bottomrule
\end{tabular}
\end{tracebox}

\begin{tracebox}[{Duplicated maximum: $s = (3, 8, 2, 8, 5)$}]
\centering
\begin{tikzpicture}
    \dparray{r}{0}{$s$}{3/adone, 8/akeep, 2/adone, 8/akeep, 5/adone}
    \dparray{rr}{-0.85}{$d$}{$-5$/compcell, 0/anscell, $-6$/compcell, 0/anscell, $-3$/compcell}
    \node[note, anchor=west, text width=4.4cm, align=left] at (5.0,-0.4) {the second 8 takes the branch \texttt{x > mx2}: $(\texttt{mx1}, \texttt{mx2})$ goes from $(8, 3)$ to $(8, 8)$; each 8 is answered with the other 8};
\end{tikzpicture}

\medskip
\small
\begin{tabular}{@{}cccccc@{}}
\toprule
\textbf{read} $x$ & 3 & 8 & 2 & 8 & 5 \\
\midrule
\texttt{mx1} & 3 & 8 & 8 & 8 & 8 \\
\texttt{mx2} & 0 & 3 & 3 & \textbf{8} & 8 \\
\bottomrule
\end{tabular}
\end{tracebox}

\subsubsection*{6. Test suite}

\begin{testbox}
\centering\small
\begin{tabular}{@{}lll@{}}
\toprule
\textbf{strengths} & \textbf{Why it is interesting} & \textbf{Expected} \\
\midrule
$5\ 5$ & $n = 2$, equal & $0\ 0$ \\
$1\ 10^9$ & $n = 2$, extreme values & $-999999999\ 999999999$ \\
$4\ 9\ 4$ & official test 4: unique maximum in the middle & $-5\ 5\ -5$ \\
$4\ 4\ 4\ 4$ & official test 5: all equal & $0\ 0\ 0\ 0$ \\
$3\ 8\ 2\ 8\ 5$ & maximum twice (distinct-second trap) & $-5\ 0\ -6\ 0\ -3$ \\
$2\ 6\ 6\ 1\ 6\ 3$ & maximum three times & $-4\ 0\ 0\ -5\ 0\ -3$ \\
$10\ 1\ 1$ & maximum first, the rest tied & $9\ -9\ -9$ \\
$t = 1$, $n = 2\cdot10^5$, all $10^9$ & largest input; output speed & $0 \dots 0$ \\
\bottomrule
\end{tabular}
\end{testbox}
Both efficient programs reproduce the official sample (five tests) and agree with the brute force on 300 random inputs of 40 tests each ($n \le 9$, strengths in $[1, 4]$, so ties are frequent).

\subsubsection*{7. Common mistakes}

\begin{warningbox}[Mistake 1 --- the distinct second maximum]
\begin{lstlisting}[style=snippet]
if (x > mx1) { mx2 = mx1; mx1 = x; }
else if (x > mx2 && x != mx1) mx2 = x;   // WRONG: (3,8,2,8,5) gives 3 for each 8
else if (x > mx2) mx2 = x;               // RIGHT: a second copy of mx1 becomes mx2
\end{lstlisting}
\end{warningbox}

\begin{warningbox}[Mistake 2 --- comparing indices instead of values]
Storing the \emph{index} of the maximum and testing \texttt{i == idx} is also correct, but then \texttt{mx2} must be the largest value at an index $\ne$ \texttt{idx}, duplicates included. Mixing the two (value test in pass 1, index test in pass 2) is a frequent source of off-by-one-copy bugs.
\end{warningbox}

\begin{warningbox}[Mistake 3 --- answering during pass 1]
Printing $s_i - \texttt{mx1}$ while reading uses the maximum of the \emph{prefix}, not of the whole array: on test 1 the first participant would get $4 - 4 = 0$. Global summaries need the full scan first.
\end{warningbox}

\begin{warningbox}[Mistake 4 --- the quadratic loop on many tests]
The brute force passes the sample and small tests and fails only on the big test. The ``sum of $n$'' guarantee does not save an $\bigO{n^2}$ algorithm: one test may have $n = 2\cdot10^5$.
\end{warningbox}

\subsubsection*{8. Takeaways}

\begin{itemize}[leftmargin=*]
    \item ``The best among the \emph{others}'' $=$ the best overall, unless you are the best; then it is the second best \emph{element}. Keep the top two with duplicates.
    \item A global summary followed by a per-position answer is the standard two-pass shape; both passes are $\bigO{n}$.
    \item Prefix and suffix aggregates answer ``everything except position $i$'' for any combinable operation, at the price of two arrays.
    \item Positive constraints ($s_i \ge 1$, $n \ge 2$) are what make the $0$ initial value safe --- check them before copying the code to another problem.
\end{itemize}

\begin{variantsbox}
\variant{Greatest element to the right}{\exeasy\ \lc{Easy}\ \varlink{Replace Elements with Greatest Element on Right Side}{LeetCode 1299}{https://leetcode.com/problems/replace-elements-with-greatest-element-on-right-side/}}
  {replace every element by the maximum of the elements \emph{to its right}; the last becomes $-1$.}
  {only the right side matters: a suffix maximum, kept as one running variable.}
  {scan right to left, write the current maximum, then fold in $a_i$; in place. $[17,18,5,4,6,1] \to [18,6,6,6,1,-1]$.}
\variant{Compare with the maximum only}{\exeasy\ \lc{Easy}\ \varlink{Kids With the Greatest Number of Candies}{LeetCode 1431}{https://leetcode.com/problems/kids-with-the-greatest-number-of-candies/}}
  {for each child, would $c_i + \mathit{extra}$ be at least the maximum of \emph{all} children?}
  {the global maximum including $i$ itself (if $i$ is the maximum, the test is trivially true).}
  {pass 1: $m_1$; pass 2: $c_i + \mathit{extra} \ge m_1$. $[2,3,5,1,3]$, extra $3 \to$ T T T F T.}
\variant{Weakest other participant}{\varnone}
  {print $s_i - \min_{j \ne i} s_j$.}
  {the two smallest elements, duplicates allowed.}
  {mirror image: answer with $\min_2$ if $s_i = \min_1$, else with $\min_1$.}
\variant{Average of the others}{\varnone}
  {print $s_i$ minus the average strength of the other $n - 1$ participants (as a fraction or real).}
  {a sum has an inverse: the total $S$ is enough, no second best is needed.}
  {$S$ in \texttt{long long} (up to $2\cdot10^{14}$); answer $s_i - (S - s_i)/(n-1)$.}
\end{variantsbox}
